\documentclass[11pt]{report}
% ============================================================
%  PREAMBOLO VERIFICA — font normale (TeX Gyre Pagella)
% ============================================================
\usepackage{fontspec}
\setmainfont{TeX Gyre Pagella}

\input{../../preambolo-verifica-core.tex}

% unicode-math va caricato dopo amsmath/amssymb/amsfonts (già caricati
% dal core), altrimenti conflitti tipo "\eth already defined".
\usepackage{unicode-math}
\setmathfont{TeX Gyre Pagella Math}


\begin{document}
\chapter{Valutazioni Verifica sulle Equazioni di II grado Incomplete}
\section*{Fila A}
\begin{description}
    \item[MELCOC] Ho il sospetto che abbia copiato da CLAMAR (es. 1 e 2 identici). Mi sembra che, a parte gli errori di base, si sia impegnata. Le equazioni di secondo grado le ha abbastanza capite. Si è confusa su molte cose che non erano eq di II grado.
    
        \textbf{Proposta: 6}
    \begin{enumerate}
        \item Molti errori (ma risultato corretto)
        \item Ha moltiplicato per $\sqrt \pi$ solo un pezzo e non tutto. Inoltre, ha dimenticato il meno.
        \item Perfetto.
        \item Manca il doppio prodotto ma a parte quello le è venuto.
        \item Ha moltiplicato per 4 solo un pezzo e non tutto.
        \item Perfetto. 
    \end{enumerate}
    \item[EDOBER] Verifica quasi perfetta (ma non ha svolto i bonus)! 
    
    \textbf{Proposta: 8½}
    \begin{enumerate}
        \item Nel primo si è perso un pochino. Il risultato è corretto, ma è un colpo di fortuna.
        \item Perfetto.
        \item Perfetto.
        \item Perfetto.
        \item Perfetto.
        \item Errorino finale di 39/36=13/2, ma a parte quello tutto perfetto.
        \item Manca
        \item Buona idea: ha riconosciuto il quadrato di un binomio ma non ha completato.
        \item Manca
    \end{enumerate}  
    \item[GUGGOL] L'impressione è che non abbia quasi fatto esercizi e abbia studiato molto poco. Però il ragazzo è bravo.
    
    \textbf{Proposta: 6}
    \begin{enumerate}
        \item Nullo.
        \item Nullo.
        \item Va bene, a parte la nozione perfetto.
        \item Nel doppio prodotto ha messo solo il prodotto senza raddoppiarlo, ma comunque ha proceduto correttamente.
        \item Non ha messo la doppia soluzione.
        \item Ha di nuovo messo il doppio prodotto come prodotto singolo, senza raddoppiare (credo sia andato a memoria). Poi ha provato a farlo con il Delta (gli veniva completa) ma ha sbagliato la formula.
        \item Manca.
        \item Nullo.
        \item Manca.
    \end{enumerate} 
    \item[CLAMAR] Ho il sospetto che abbia fatto copiare MELCOC (esercizi 1 e 2 identici).
    
    \textbf{Proposta: 6}
    \begin{enumerate}
        \item Stessi errori di MELCOC.
        \item Stessi errori di MELCOC.
        \item Perfetta (come MELCOC, anche se procedimento leggermente diverso).
        \item Perfetto.
        \item Quando aveva $\frac{x^2}4=12$ ha diviso per 4 invece di moltiplicare per 4, ma a parte quello ha fatto tutto bene.
        \item Stesso curiosissimo errore di Cochis ma in realtà il risultato è corretto.
        \item Nullo.
        \item Ha riconosciuto il quadrato del binomio ma non ha messo le due soluzioni e ne ha messa solo una.
        \item Manca.
    \end{enumerate}
\end{description}
\section*{Fila B}
\begin{description}
    \item[GINVAL] In generale mi sembra che abbia studiato e abbia anche compreso. A volte ha un buon intuito ed è stata anche precisa con i calcoli.
    
    \textbf{Proposta: 8-}
    \begin{enumerate}
        \item Perfetto.
        \item Nullo (è andata a memoria e ha sbagliato).
        \item Perfetto.
        \item Perfetto.
        \item Si è dimenticata il più o meno alla fine.
        \item Si è solo dimenticata un meno, ma per il resto tutto perfetto.
        \item Molto bene, è andata a intuito e ha trovato, solo guardandola, che 0=0 e quindi t=1.
        \item Ha utilizzato la formula con il Delta (quasi giusta però... solo ha sbagliato il segno davanti a b alla fine). Nullo.
        \item Nullo.
    \end{enumerate}
\item[RICRAO] Secondo me ha studiato poco. Inoltre, gli ho tolto un punto perché parlava durante la verifica.

\textbf{Proposta: 5-1=4.} (Il 5 è generoso).
\begin{enumerate}
    \item Tanta confusione. Non ha ragionato sui segni. Nullo.
    \item Non ha neanche messo il quadrato su x... nullo.
    \item Perfetto.
    \item Non ha saputo fare il quadrato del binomio. Nullo.
    \item Ha fatto bene il prodotto notevole ma poi quando è arrivato a $t^2=4$ ha concluso t=4! Nullo.
    \item Molti pasticci. Nullo.
    \item Manca.
    \item Nullo.
    \item Manca.
\end{enumerate}

\item[PIEBAI] Secondo me ha studiato un po', credo si sia abbastanza impegnato. Ha provato a fare praticamente tutto, è venuto alla correzione della simulazione, in classe chiedeva di fare esercizi.

\textbf{Proposta: 6½/7}
\begin{enumerate}
    \item Troppi errori. Nullo.
    \item Perfetto.
    \item Nullo.
    \item Perfetto.
    \item Perfetto.
    \item Alcuni errori di distrazione ma il succo c'è.
    \item Nullo.
    \item Nullo.
    \item Manca.
\end{enumerate}
\item[ARIDAN] Ha lavorato benissimo. Anche in classe.

\textbf{Proposta: 10}
\begin{enumerate}
    \item Perfetto.
    \item Perfetto.
    \item Perfetto.
    \item Perfetto.
    \item Perfetto.
    \item Perfetto.
    \item Perfetto.
    \item Perfetto.
    \item Manca.
\end{enumerate}

\item[NINTRO] Fa tanti errori, non si impegna molto e secondo me non sapeva le spurie. Però ha imparato le pure e le monomie.

\textbf{Proposta: 6/6½}
\begin{enumerate}
    \item Perfetto.
    \item Nullo.
    \item Perfetto (sospetto un minimo che abbia copiato da GINVAL).
    \item Si è bloccata dopo che è arrivata all'eq spuria.
    \item Ha fatto un errore di calcolo irrilevante ma l'esercizio è perfetto.
    \item Molti errorini. Giunge a un'equazione completa, che quindi non può risolvere.
    \item Manca.
    \item Manca. 
    \item Manca.
\end{enumerate}

\item[GUICER] Per sua stessa ammissione non ha fatto i compiti. Non ha studiato palesemente nulla.

\textbf{Proposta: 4½/5}
\begin{enumerate}
    \item Nullo.
    \item Ha mancato le due soluzioni: ne ha scritta solo una.
    \item Giusto. Poi l'ha rifatto in un altro modo che invece è delirante ma vabbè.
    \item Ha sbagliato il quadrato di un binomio e non sa risolvere una spuria.
    \item Non ha messo la doppia soluzione.
    \item Ha sbagliato il meno davanti alla parentesi, quindi gli veniva un'equazione completa. Poi per risolverla ha diviso per x.
    \item Ha scritto $(t-1)^7=t^7-1$ e $(t-1)^5 = t^5-1$.
    \item Diviso per t, nullo.
    \item Vari errori. Nullo.
\end{enumerate}

\end{description}

\section*{Fila C}

\begin{description}

\item[DYLPAS] Secondo me si è impegnato fin dall'inizio e ha fatto molto bene. Ha capito come fare tutti gli esercizi. Ha fatto solo qualche errore di distrazione.

\textbf{Proposta: 7½/8}
\begin{enumerate}
    \item Perfetto.
    \item Perfetto.
    \item Perfetto.
    \item Un po' di errori di distrazione, non ha riconosciuto che era impossibile.
    \item Un errorino di distrazione ma il succo c'è.
    \item Un paio di errorini di distrazione ma il succo c'è.
    \item Ha sviluppato il prodotto,e chiaramente non ha funzionato.
\end{enumerate}

\item[CARANT] Non è stata attenta, non mi sembra si sia impegnata. Mi ha anche inviato gli schemi sabato di mezzanotte e mezza, quelli schemi non andavano bene e lei si è rifiutata di togliere gli esempi. Quantomeno, durante la verifica, ha avuto un buon atteggiamento. Non credo abbia imparato a risolvere nessun tipo di equazione.

\textbf{Proposta: 4½}

\begin{enumerate}
    \item Perfetto (gliel'avevo un po' suggerito io).
    \item Manca il più o meno davanti alle soluzioni e poi ha fatto spuntare un'altra radice che boh.
    \item Nullo.
    \item Ha fatto bene fino a $z^2=-1$, poi non ha saputo concludere.
    \item Sbagliato lo svolgimento del quadrato di un binomio, sbagliato a sommare i termini simili, si è fermata quando è arrivata a x²-2=0, dimostrando di non saperla risolvere.
    \item Ha scritto $-(x-1)^2=(x+1)^2$ e poi ha svolto il quadrato del binomio pure male. Poi è arrivata a x²+x=0 e non ha saputo procedere.
\end{enumerate}

\item[VALCAS] \textbf{Non si è portata gli schemi.} Non ha fatto niente di corretto: a volte isolava x² ma poi non sapeva continuare. Sulle equazioni di seocndo grado non ha imparato nulla. 

\textbf{Proposta: 4}

\item[PIECOL] Qualcosina ha fatto ma non molto. Ha fatto un po' di errori di distrazione, altri più concettuali. Comunque sufficiente.

\textbf{Proposta: 7-}

\begin{enumerate}
    \item Perfetto (anche se l'ho dovuto scongiurare di scrivere x=0 e non lasciarlo come $x=\pm\sqrt{\frac 02}$)
    \item Ha messo $y=\pm\sqrt 4$ ma poi ha scritto $y=2$. Glielo conterei giusto (o quasi).
    \item Perfetto.
    \item Ha scritto 4/4=0, quindi si è sfanculata l'eq.
    \item A parte $x-1=0 \implies x=-1$ ha fatto tutto perfetto.
    \item Ha sbagliato il quadrato del binomio (non ha messo il meno). poi si è inventato che $-x(x+2)=2$ si risolve come le spurie: $x=0 \vee x+2=2$.
    \item Manca.
    \item Manca.
    \item Manca.
\end{enumerate}

\end{description}

\end{document}